\label{sec:diagnose}

We diagnose two capabilities within the interaction loop
$B_t\rightarrow P_t\rightarrow e_{t+1}\rightarrow B_{t+1}$:
forming a task-relevant partner judgment ($B$) and using it to plan an
interaction ($P$). The diagnosis separates two questions: \emph{what dependencies
does the task impose?} and \emph{can the model perform the computations required
at each interface?} An unreliable planner can fail to use a correct judgment,
and an unreliable updater can fail to use informative evidence. Consequently,
the final model output alone cannot distinguish an absent dependency from a
failure to exploit it. We first certify both directions with an oracle, then
measure model errors under controlled judgments and evidence. The effect of
replacing one model component is a supplementary observation, not the criterion
for whether a task dependency exists.

\subsection{Native endgames and paired interventions}

\paragraph{Controlled positions.}
We freeze legal public histories in three-player BENAC-P games and measure at
most two actual ego decisions, including any intervening response decision.
The original proposal schedule, full legal action set, independent preference
catalogues, and irreversible commitments are retained. Partners actively
propose and respond using the history-aware deterministic policy in
Section~\ref{sec:benac-p}. In each position, one partner--goal preference is
queried; all other preference entries are supplied, and uncertainty is over
\textsc{Want}, \textsc{Neutral}, and \textsc{Avoid}. Exact search supplies
reference judgments and terminal-utility action values, without exposing the
realized hidden preference to ego.

We use two oracle-selected cohorts of six positions each, sharing two
positions: ten positions from ten source games in total. The $B\rightarrow P$
cohort requires an action that is optimal under one supported partner type but
strictly suboptimal under another. Common optimal actions are allowed and
remain correct. The $P\rightarrow B$ cohort requires two legal first actions
that induce different expected posterior uncertainty and leave a next ego
decision on every response branch. These forced actions need not be
utility-optimal; the reference action used for planning regret always maximizes
terminal utility over all legal actions. Selection uses game mechanics and
oracle values, not model errors or the sign of a repair effect.

\paragraph{A concrete negotiation.}
Consider \texttt{native-30040-g6-d2}, where ego is $P_1$, all three players
have already committed to $A_0$, and the remaining proposers are $P_1$ then
$P_0$. The queried preference is $P_2$'s attitude toward
$G_6=(P_0.A_1\land P_2.A_1)$. Ego can offer $O_1$ alone or the menu
$\{O_1,O_2\}$, with new commitments
\[
 O_1:\;P_1\text{ adds nothing},\;P_2\text{ adds }A_1;
 \qquad
 O_2:\;P_1\text{ adds }A_1,\;P_2\text{ adds }A_1.
\]
Under the declared oracle policy and this game's other preferences, all three
$P_2$ types accept $O_1$ alone, leaving the queried preference unresolved.
With the menu, however, \textsc{Avoid} chooses $O_1$, whereas \textsc{Want}
and \textsc{Neutral} choose $O_2$. Thus the choice supports either
$\{\textsc{Avoid}\}$ or $\{\textsc{Want},\textsc{Neutral}\}$, while immediately
binding the selected commitments. Next, $P_0$ actively proposes adding its own
$A_1$ without further commitments from $P_1$, who must accept or reject.
This is evidence acquisition followed by a native decision, not an explicit
query action. In the observed $O_1$ menu branch, the model reports
$\{\textsc{Want},\textsc{Neutral}\}$ despite the response supporting
$\{\textsc{Avoid}\}$. Appendix~\ref{app:diagnose-example} gives the preferences
and oracle payoff comparison underlying this example.

\paragraph{Separate capability measurements.}
For $B$, the model receives the initial possibilities and public history and
returns the set of partner preferences still compatible with that history.
We score exact set agreement, allowing a correct answer to remain uncertain.
For $P$, we supply the evidence-supported reference judgment and ask the model
to select a legal action. Regret is the difference from the optimal expected
terminal utility, with reference continuation after the chosen action; all
optimal ties receive zero regret. The model outputs brief reasoning followed
by a semantic judgment or action index, never probabilities, $Q$ values, or a
plan variable. Since $P$ also receives history, this is planning assisted by a
correct judgment, rather than isolation of an internal planning module.

\paragraph{Certifying dependencies and probing their use.}
For $B\rightarrow P$, exact search holds the physical state, public history
and legal actions fixed and compares action values under singleton judgments
for two still-compatible partner types. An action optimal under one judgment
must be strictly suboptimal under the other. This certifies that partner
understanding affects action evaluation; shared optimal actions remain valid.
For $P\rightarrow B$, we force two legal first actions from the same root,
enumerate subsequent partner actions up to the next ego decision, and compute
the evidence-supported judgments and their expected uncertainty. This certifies
that interaction choice affects what can subsequently be inferred. Both are
oracle selection criteria, established independently of model performance.

We then ask the model to form judgments from those histories and to plan with
the evidence-supported reference judgment. These tests locate errors in $B$
and in reference-judgment-assisted $P$, respectively. As a supplementary
$B\rightarrow P$ probe, we replace the model's explicit judgment with the
reference judgment in otherwise identical fresh planning contexts. The
replacement supplies only what the evidence supports, not the realized hidden
type. Action changes and signed regret differences measure the model's response
to this input intervention; neither a positive repair benefit nor a gain in
accuracy between evidence arms is required to certify the task dependencies.
Because actions also change commitments, utility contrasts are not interpreted
as pure information effects.

\subsection{Certified dependencies and model errors}

We evaluate Qwen3-4B-Instruct-2507 at temperature zero. All 127 diagnostic tasks
returned valid submissions, with no truncation and a mean of 230.8 completion
tokens. The separate empty-history controls scored 3/4: a known
\textsc{Want} preference was incorrectly expanded to all three possibilities.
We retain this error and therefore do not attribute all subsequent belief
errors exclusively to strategic inference. The cohorts are development
diagnostics, not an untouched confirmation set. Table~\ref{tab:diagnose-main}
separates oracle-certified dependencies from model measurements; each direction
uses its own six-game cohort. Appendix~\ref{app:diagnose} gives
selection, execution and aggregation details.

\begin{table}[t]
\caption{Two levels of diagnosis. Panel A reports exact task properties required
for selection, not model success rates. Panel B measures model performance at
the corresponding interfaces. Each direction uses six source games, with two
shared across directions. Intervals are descriptive 95\% game-bootstrap
intervals; regret uses additive goal-utility units.}
\label{tab:diagnose-main}
\centering
\small
\begin{tabular}{lrr}
\hline
Diagnostic quantity & Result & 95\% interval \\
\hline
\multicolumn{3}{l}{\emph{A. Oracle-certified task dependencies}} \\
$B\rightarrow P$: judgment changes action optimality & 6/6 cases & --- \\
$P\rightarrow B$: action changes inferable preferences & 6/6 cases & --- \\
Posterior entropy: lower / higher information (bits) & 1.585 / 0.667 & --- \\
\hline
\multicolumn{3}{l}{\emph{B. Model errors at the two interfaces}} \\
Root B exact, judgment-sensitive cohort & 16.7\% & [0.0, 50.0]\% \\
P regret with oracle judgment, judgment-sensitive cohort & 0.444 & [0.167, 0.722] \\
P regret with oracle judgment, action--evidence cohort & 1.000 & [0.556, 1.556] \\
B exact after forced lower-information action & 16.7\% & [0.0, 50.0]\% \\
B exact after forced higher-information action & 22.2\% & [5.6, 44.4]\% \\
\hline
\end{tabular}
\end{table}

\paragraph{The task couples partner understanding and interaction choice.}
Panel A makes the two dependencies explicit. In every judgment-sensitive case,
changing the assumed partner type changes the optimality of at least one legal
action. In every action--evidence case, changing the forced interaction changes
the set of partner preferences that subsequent behavior supports. The latter
reduces expected oracle uncertainty by 0.918 bits for each selected pair.
Thus, both dependencies exist even before considering whether the model can
compute an appropriate judgment or action. These are properties of the
selected endgames, not estimates of their prevalence in random games.

\paragraph{The model makes errors at both interfaces.}
Panel B tests the model against those oracle-defined requirements. At the
judgment-sensitive roots, only one of six reported partner judgments is exact.
When the reference judgment is supplied, planning regret remains positive:
0.444 in the judgment-sensitive cohort and 1.000 in the action--evidence
cohort. The planning test therefore does not require the model first to infer
the correct judgment unaided. Conversely, the forced-action tests provide the
corresponding evidence independently of the model's first-action choice.
Judgment accuracy is 16.7\% under the lower-information action and 22.2\% under
the higher-information action. A correct unresolved set receives full credit;
these errors concern agreement with what the evidence actually supports, not
failure to identify a type that remains ambiguous.

\paragraph{Joint responses are supplementary, not a dependency test.}
Replacing model judgments changes three of six root actions in the
judgment-sensitive cohort, with two utility-neutral changes and one decrease
in utility. The mean regret reduction is $-0.111$; the paired difference in
B accuracy between the forced evidence arms is 5.6 percentage points, with an
interval spanning zero. Appendix~\ref{app:diagnose} reports these signed effects.
They describe how the evaluated model responds to the interventions. An output
that remains wrong after receiving a correct judgment or more informative
evidence does not erase the independently certified task dependency.

The diagnosis therefore identifies errors in both partner judgment and
partner-conditioned planning within tasks that couple the two computations.
It separates the existence of those dependencies from the model's ability to
use them. This motivates training $B$ and $P$ as separate targets and then
evaluating their composition and transfer. It does not identify internal neural
modules or establish that the two errors causally amplify one another.
